% =============================================================================
% Preambulo do documento "Como o Morphe funciona" (docs/como-funciona/). E o
% mesmo dos guias e roteiros: caixas Por que / Faca / Confira / Cuidado e um
% diagrama por parte. Compila no Overleaf com pdfLaTeX (suba os dois .tex)
% ou com Tectonic.
% =============================================================================

\usepackage[utf8]{inputenc}
\usepackage[T1]{fontenc}
\usepackage{lmodern}
\usepackage[brazilian]{babel}
\usepackage[a4paper,margin=2.2cm]{geometry}
\usepackage{microtype}
\setlength{\emergencystretch}{3em}
\setlength{\parindent}{0pt}
\setlength{\parskip}{4pt}
\usepackage{xcolor}
\usepackage{booktabs}
\usepackage{tabularx}
\usepackage{array}
\usepackage{enumitem}
\setlist{itemsep=1pt,topsep=3pt}
\usepackage[most]{tcolorbox}
\usepackage{listings}
\usepackage{tikz}
\usetikzlibrary{positioning,arrows.meta,fit,calc,shapes.geometric,backgrounds}
\usepackage[hidelinks]{hyperref}

% No pdfLaTeX (Overleaf), os espacos do codigo viram so distancia no PDF, e
% copiar "git clone -b" cola "git clone-b". Isto grava um espaco de verdade
% entre as palavras; o XeTeX (Tectonic) ja faz isso sozinho.
\ifdefined\pdfinterwordspaceon
  \input{glyphtounicode}\pdfgentounicode=1
  \pdfinterwordspaceon
\fi

% --- cores: azul = computador/aplicativo, verde = placa (HPS), laranja = FPGA ---
\definecolor{cli}{HTML}{2563EB}
\definecolor{hps}{HTML}{15803D}
\definecolor{fpga}{HTML}{C2410C}
\definecolor{fundo}{HTML}{F3F4F6}
\definecolor{cinza}{HTML}{6B7280}
\definecolor{verm}{HTML}{DC2626}
\definecolor{git}{HTML}{7C3AED}

\newcommand{\arq}[1]{{\upshape\ttfamily\detokenize{#1}}}
\newcommand{\opt}[1]{{\upshape\ttfamily -{}-\detokenize{#1}}}

% --- caixas ---
\newtcolorbox{porque}{colback=fpga!4,colframe=fpga!70!black,boxrule=0.4pt,arc=2pt,
  title=Por quê,fonttitle=\bfseries\small,fontupper=\small,left=4pt,right=4pt}
\newtcolorbox{faca}{breakable,colback=white,colframe=black!60,boxrule=0.4pt,arc=2pt,
  title=Faça,fonttitle=\bfseries\small,fontupper=\small,left=4pt,right=4pt}
\newtcolorbox{confira}{breakable,colback=green!5,colframe=green!45!black,boxrule=0.4pt,arc=2pt,
  title=Confira,fonttitle=\bfseries\small,fontupper=\small,left=4pt,right=4pt}
\newtcolorbox{cuidado}{breakable,colback=orange!7,colframe=orange!75!black,boxrule=0.4pt,arc=2pt,
  title=Cuidado,fonttitle=\bfseries\small,fontupper=\small,left=4pt,right=4pt}
\newtcolorbox{pronto}{colback=blue!4,colframe=blue!55!black,boxrule=0.4pt,arc=2pt,
  title=Já está pronto,fonttitle=\bfseries\small,fontupper=\small,left=4pt,right=4pt}

% --- codigo: um comando por bloco, sem barra invertida, para copiar e colar ---
\lstset{
  basicstyle=\ttfamily\footnotesize, backgroundcolor=\color{fundo},
  commentstyle=\color{cinza}, frame=single, framerule=0pt, framesep=3pt,
  xleftmargin=4pt, xrightmargin=4pt, breaklines=true, columns=fullflexible,
  keepspaces=true, showstringspaces=false, upquote=true,
  aboveskip=4pt, belowskip=4pt,
  literate={á}{{\'a}}1 {é}{{\'e}}1 {í}{{\'i}}1 {ó}{{\'o}}1 {ú}{{\'u}}1
           {ã}{{\~a}}1 {õ}{{\~o}}1 {ç}{{\c c}}1 {ê}{{\^e}}1 {ô}{{\^o}}1
}

% --- diagramas ---
\tikzset{
  bloco/.style={draw,rounded corners=2pt,align=center,font=\small,
                minimum height=0.9cm,inner sep=4pt},
  cli/.style={bloco,fill=cli!10,draw=cli},
  hps/.style={bloco,fill=hps!10,draw=hps},
  fpga/.style={bloco,fill=fpga!10,draw=fpga},
  gitb/.style={bloco,fill=git!10,draw=git},
  neutro/.style={bloco,fill=black!4,draw=black!50},
  seta/.style={-{Stealth[length=2mm]},thick},
  rot/.style={font=\scriptsize,midway,align=center},
  passo/.style={rounded corners=2pt,fill=black,text=white,font=\bfseries\scriptsize,
                inner sep=2pt,minimum size=5mm},
}
\newcommand{\passo}[1]{\tikz[baseline=-0.6ex]\node[passo]{#1};}
\newcommand{\etapa}[2]{\subsection*{\passo{#1}\ \ #2}\addcontentsline{toc}{subsection}{#1\ \ #2}}
